\section{Resumen}

Los mercados cripto híbridos combinan mecanismos centralizados de negociación con
creadores automáticos de mercado descentralizados (AMM). El arbitraje puede reducir
discrepancias entre precios, mientras que restricciones de ejecución, decisiones
heterogéneas de negociación y cambios endógenos de liquidez pueden modificar la
dinámica de ajuste. Se propone un marco metodológico con dos niveles: (i) un mapa
discreto reducido cuyas variables son observables en el simulador; y (ii) un micromundo
estocástico con un precio CEX de referencia, un AMM de producto constante y cuatro
tipos de agentes. Simulaciones sintéticas preliminares verifican las direcciones
mecánicas esperadas: el arbitraje acota el error de seguimiento, menor capacidad
prolonga la recuperación y el retiro adaptativo de liquidez amplifica el impacto de una
venta fija. Esas pruebas no validan el modelo frente a mercados reales ni convierten los
coeficientes reducidos en parámetros empíricos. La investigación contrastará el mapa
reducido con el ABM mediante escenarios controlados, sensibilidad y resultados de
desacuerdo.

\textbf{Palabras clave:} seguimiento CEX/DEX; AMM; arbitraje; liquidez adaptativa;
simulación multiagente; mapa discreto; reproducibilidad.
